\section{Ethereum as a State-Transition System}
Ethereum is a blockchain that can be viewed as a deterministic replicated
state machine~\cite{ethereum22,buterin2014ethereum}.
A transaction is
one execution which consists of a sequence of state transitions.
Let $\q$ denote the global blockchain state
immediately before transaction execution,
and let $\q'$ denote the resulting state after
execution. For a transaction $\tau$, the execution
relation is written as $\q \xrightarrow{\tau} \q'$ when
execution succeeds, or as $\q \xrightarrow{\tau} \bot$ when
execution reverts.

At the consensus layer, blocks impose a total order over
transactions. At the execution layer, each transaction
is interpreted as a sequence of Ethereum Virtual Machine (EVM)
state transitions
constrained by gas, call semantics, and exception propagation.
This deterministic transition view is the basis
of empirical transaction replays used in this thesis to
evaluate the effectiveness of runtime defenses.
This state-transition perspective also provides
the data collection used later for invariant generation
 and runtime guard deployment.

Ethereum maintains two account classes:
\begin{itemize}
    \item \textbf{Externally owned accounts (EOAs):}
    controlled by cryptographic keys and capable of
    initiating top-level transactions.
    \item \textbf{Contract accounts:} controlled
    by deployed bytecode and activated via message calls.
\end{itemize}

Each account is associated with an address,
a nonce, a native-token balance, and storage-relevant
metadata. Contract accounts additionally
contain executable bytecode and persistent key-value
storage. High-level languages such as
Solidity~\cite{solidity} and Vyper~\cite{vyper} compile to
EVM bytecode, and this bytecode is the only
representation executed on-chain. Figure
~\ref{fig:ERC20} shows an example of a Solidity contract.

% This distinction is important for security
% reasoning: source-level patterns are informative,
% but exploitability is ultimately determined
% by bytecode-level behavior under EVM semantics.

\input{solidity-highlighting.tex}
\begin{figure}[t]
    \centering
    \begin{minipage}[h]{0.9\linewidth}
    \lstinputlisting{codes/ERC20.sol}
    \end{minipage}
    \caption{An excerpt of ERC20 smart contract.}
    \label{fig:ERC20}
    \vspace{-2mm}
\end{figure}


\subsection{EVM Semantics and Execution}

EVM refers to the Ethereum Virtual Machine, which is the runtime
environment for executing smart contract bytecode.
An EVM execution frame is characterized by program counter,
stack, memory, gas counter, and an execution context
(caller, callee, call value, calldata, and return-data buffers).
Execution proceeds by repeatedly decoding
and applying opcodes until a halting condition is reached.
EVM consists a list of
low-level opcodes, such as \texttt{ADD}, \texttt{SLOAD},
which manipulate the execution frame and global state,
and perform state transitions.

Semantically, opcode execution induces a small-step
relation over machine configurations. Multi-contract
transactions form a call tree through
message calls (e.g., \texttt{CALL},
\texttt{DELEGATECALL}, \texttt{STATICCALL}),
where child-frame outcomes feed back into parent
frames via return values or propagated exceptions.
This call-tree structure explains why cross-contract
analysis is necessary even when auditing one protocol.

% The remaining subsections characterize the execution
% dimensions most relevant to security behavior.

\subsection{EVM Storage Classes}
The EVM exposes storage classes with distinct
persistence properties:
\begin{itemize}
    \item \textbf{Stack and memory:}
    ephemeral to the current call frame.
    \item \textbf{Persistent storage:}
    contract-scoped and transaction-independent;
    updates survive across blocks, accessed via
    opcodes such as \texttt{SLOAD} and \texttt{SSTORE}.
    \item \textbf{Transient storage:}
    transaction-scoped scratch space accessible
    by recent introduced EVM opcodes
    (e.g., \texttt{TLOAD}/\texttt{TSTORE});
    values are reset at transaction boundaries.
\end{itemize}
These distinctions are not merely implementation
details: they also shape both attack feasibility
and defense cost models which will be explored in later chapters.

\subsection{Gas and Transaction Atomicity}
Every EVM opcode consumes gas. Gas is a measure of
computational effort and is paid by the transaction sender,
in units of the native token (Ether).
 A transaction specifies
a gas limit and terminates successfully only if
cumulative consumption remains within this budget
and no uncaught exception occurs.
There is an upper limit on
per-block gas consumption, which constrains the
feasibility of complex transactions and deep call trees.
Out-of-gas errors, invalid opcodes, and
explicit \texttt{REVERT} induce exceptional
control flow; in particular, state changes
in the failing frame (and descendants) are
rolled back according to call semantics.
Gas therefore can act simultaneously as an economic
mechanism and a semantic guardrail.

At transaction level, Ethereum execution
is atomic: either all committed effects of
the top-level frame are recorded in $\q'$,
or the transaction is rejected and the pre-state $\q$
is preserved. Within a transaction,
nested calls can fail locally and be handled by parent
code, but unhandled failures propagate and may
invalidate larger portions of execution until either
the transaction fails or a successful error handle is reached.

This atomic behavior is critical for
reasoning about DeFi attacks, invariant
enforcement, and control-flow guards, where success
often depends on subtle interactions between local
checks and global rollback semantics.
Consequently, any analysis framework that
reasons only about local function semantics
is likely to be incomplete for compositional
attack scenarios.

\subsection{Transactions and Inter-Contract Composition}
A transaction contains sender, recipient, nonce,
value, calldata, and gas parameters. If the
recipient has no code, the effect is Ether
transfer; otherwise the recipient bytecode
is invoked with calldata-decoded arguments.

In modern DeFi smart contract systems, a single top-level
transaction can trigger complex inter-contract
compositions: exchanges, lending pools,
vaults, and oracle updates may all be traversed
before termination. Consequently, transaction
behavior is determined jointly by local code and global
execution context induced by prior calls and shared
state dependencies.

This observation directly motivates the later
technical chapters: each method in this thesis is
designed to reason about traces, rather than
isolated function bodies.

\subsection{DeFi Protocol}
Decentralized finance (DeFi) comprises
financial protocols implemented as composable
smart contracts~\cite{wust2018you}. Core primitives
include decentralized exchanges,
lending/borrowing markets, stablecoins, and oracle services.
A distinctive property is \emph{composability}:
protocol endpoints can be chained atomically
in one transaction.
This composability enables efficient financial
workflows but also creates many attack
surfaces. Flash loans, for example, permit
large uncollateralized borrowing if and only
if repayment occurs within the same transaction;
this feature amplifies adversarial capital
and directly interacts with EVM atomicity and
rollback semantics.
In summary, EVM semantics and DeFi composability
are inseparable in security analysis.
